%%%%%%%%%%%%%%%%%%%%%%%%%%%%%%%%%%%%%%%%%%%%%%%%%%%%%%%%%%%%%%%%%%%%%%%%%%%
%% NIER @ ASE 2026 -- Anonymous double-blind submission
%% "Distributional Program Analysis: Treating Large Language Models
%%  as Program Samplers"
%%
%% Build: pdflatex main && bibtex main && pdflatex main && pdflatex main
%%%%%%%%%%%%%%%%%%%%%%%%%%%%%%%%%%%%%%%%%%%%%%%%%%%%%%%%%%%%%%%%%%%%%%%%%%%

%TC:macro \cite [option:text,text]
%TC:macro \citep [option:text,text]
%TC:macro \citet [option:text,text]
%TC:envir table 0 1
%TC:envir table* 0 1
%TC:envir tabular [ignore] word
%TC:envir displaymath 0 word
%TC:envir math 0 word
%TC:envir comment 0 0

\PassOptionsToPackage{disable}{microtype}
\documentclass[sigconf,review,anonymous,nonacm]{acmart}

\AtBeginDocument{%
  \providecommand\BibTeX{{\normalfont B\kern-.05em{\scshape i\kern-.025em b}\kern-.08em\TeX}}}

\setcopyright{none}
\acmConference[ASE '26]{40th IEEE/ACM International Conference on Automated Software Engineering --- NIER Track}{October 12--16, 2026}{Munich, Germany}
\acmISBN{}
\acmDOI{}
\settopmatter{printacmref=false}
\renewcommand\footnotetextcopyrightpermission[1]{}
\pagestyle{plain}

%% --- Packages ----------------------------------------------------------
\usepackage{listings}
\usepackage{xcolor}
\usepackage{booktabs}
\usepackage{balance}
\usepackage{amsmath,amssymb}

\definecolor{codebg}{rgb}{0.97,0.97,0.93}
\definecolor{codecomment}{rgb}{0.40,0.45,0.55}
\definecolor{codekw}{rgb}{0.10,0.30,0.65}
\lstdefinestyle{code}{
  backgroundcolor=\color{codebg},
  basicstyle=\ttfamily\scriptsize,
  commentstyle=\color{codecomment}\itshape,
  keywordstyle=\color{codekw}\bfseries,
  numbers=left,
  numberstyle=\tiny\color{codecomment},
  numbersep=4pt,
  frame=single,
  framesep=3pt,
  rulecolor=\color{codecomment},
  language=Python,
  showstringspaces=false,
  breaklines=true,
  keepspaces=true,
  columns=fullflexible,
  aboveskip=4pt,
  belowskip=4pt,
}
\lstset{style=code}

\newcommand{\Pr}{\mathbb{P}}
\renewcommand{\Pr}{\mathbb{P}} % override only if needed
\newcommand{\samp}{\sim}

%% --- Front matter -------------------------------------------------------
\begin{document}

\title{Distributional Program Analysis:\\
Treating Large Language Models as Program Samplers}

\author{Anonymous Author(s)}
\affiliation{%
  \institution{Submission to NIER @ ASE 2026}
  \country{}}
\email{anonymous@example.org}

\renewcommand{\shortauthors}{Anonymous}

\begin{abstract}
A specification under-determines a program. A large language model (LLM),
prompted with that specification, induces a distribution over programs
that purport to satisfy it. Existing pipelines that pair static analysis
with LLM-generated code throw away the most informative artifact this
distribution provides: the \emph{variation across samples}. We propose
\textbf{distributional program analysis} (DPA): run an analyzer not on
one program but across many LLM samples for the same task, and treat the
analyzer's outputs as posterior evidence about the latent intended
program. Robust invariants---those holding in a high fraction of
samples---are likely intended; fragile invariants are likely artifacts of
decoding noise. The analyzer becomes a Bayesian critic; the LLM becomes
a sampler. We formalize survival rate as a posterior estimator, sketch a
prototype, and present a pilot on HumanEval+ showing that high-survival
invariants align with reference-solution invariants in roughly four out
of five cases, and that filtering samples by posterior consistency
improves pass@1 at fixed compute.
\end{abstract}

\begin{CCSXML}
<ccs2012>
 <concept>
  <concept_id>10011007.10011006.10011008</concept_id>
  <concept_desc>Software and its engineering~General programming languages</concept_desc>
  <concept_significance>300</concept_significance>
 </concept>
 <concept>
  <concept_id>10011007.10011074.10011099</concept_id>
  <concept_desc>Software and its engineering~Software verification and validation</concept_desc>
  <concept_significance>500</concept_significance>
 </concept>
 <concept>
  <concept_id>10010147.10010178.10010179</concept_id>
  <concept_desc>Computing methodologies~Natural language processing</concept_desc>
  <concept_significance>300</concept_significance>
 </concept>
</ccs2012>
\end{CCSXML}

\ccsdesc[500]{Software and its engineering~Software verification and validation}
\ccsdesc[300]{Software and its engineering~Automated static analysis}
\ccsdesc[300]{Computing methodologies~Natural language processing}

\keywords{program analysis, large language models, abstract interpretation,
invariant inference, Bayesian inference, code generation}

\maketitle

%% --- 1. Introduction ----------------------------------------------------
\section{Introduction}
\label{sec:intro}

A specification rarely picks out a single program. Two correct
implementations of \texttt{is\_prime} can differ in branching, data
representation, and the exact set of inputs they handle gracefully. When
that specification is fed to a large language model (LLM), the model
samples from a \emph{distribution} of plausible programs whose support
includes both the intended implementation and a long tail of
near-miss variants. Existing pipelines that follow LLM code generation
with static analysis treat this distribution as a nuisance: they pick
one sample, analyze it, accept or reject. This paper argues that the
distribution is the most informative signal in the entire pipeline---and
that current practice systematically discards it.

We make a single conceptual move. Treat the LLM as a \emph{prior} over
programs that satisfy a specification $\phi$. Treat a static analyzer as
a \emph{critic} that produces, for each sample, a set of candidate
invariants (ranges, nullability, aliasing, monotonicity, post-conditions).
Aggregate per-invariant \emph{survival rates} across $N$ samples. An
invariant that holds in 95\% of samples is likely a property of the
\emph{latent intended program}; one that holds in 30\% is likely a
property of decoding noise. The analyzer is no longer a per-artifact
guard; it is a Bayesian critic that accumulates evidence about the
intent the prompt under-specified.

We call this idea \textbf{distributional program analysis} (DPA). The
contribution of this paper is the framing, a precise definition of
posterior invariant survival, a prototype that instantiates DPA on top
of an off-the-shelf abstract interpreter, and a pilot on HumanEval+
\cite{liu2023humanevalplus} showing that (i) survival rate correlates
strongly with whether an invariant appears in the reference solution
($r \approx 0.78$ in our pilot) and (ii) filtering candidate samples by
agreement with high-survival invariants improves pass@1 at \emph{fixed}
inference compute relative to majority voting.

\smallskip\noindent\textbf{Contributions.}
\begin{itemize}
\item A new framing in which an LLM is the prior, an analyzer is the
critic, and \emph{invariant survival} is a posterior estimator over the
intended program (\S\ref{sec:approach}).
\item A formal definition of survival-based posterior weighting and a
prototype on Python with an abstract-interpretation back-end
(\S\ref{sec:approach}).
\item A pilot on HumanEval+ (\S\ref{sec:pilot}) showing the survival
estimator aligns with ground-truth invariants and improves pass@1 at
matched compute.
\item An evaluation plan (\S\ref{sec:plan}) with three sharp questions
about generality, contamination, and the substitutability of analysis
for self-consistency~\cite{wang2023selfconsistency}.
\end{itemize}

%% --- 2. Motivating example ----------------------------------------------
\section{A Motivating Example}
\label{sec:motivation}

Consider the prompt: \emph{``Implement \texttt{rolling\_max(xs, k)} that
returns the maximum value in each length-$k$ window of \texttt{xs}.''}
We sample $N=20$ implementations from a frontier LLM. Of those, 18 pass
the canonical doctest. Listing~\ref{lst:variants} shows three
representative samples.

\begin{lstlisting}[caption={Three samples for the same prompt. All pass the
canonical doctest.},label={lst:variants}]
# Sample A
def rolling_max(xs, k):
    return [max(xs[i:i+k]) for i in range(len(xs)-k+1)]

# Sample B
def rolling_max(xs, k):
    out, dq = [], []
    for i, x in enumerate(xs):
        while dq and dq[0] <= i - k: dq.pop(0)
        while dq and xs[dq[-1]] < x:  dq.pop()
        dq.append(i)
        if i >= k - 1: out.append(xs[dq[0]])
    return out

# Sample C
def rolling_max(xs, k):
    if k <= 0: return []
    return [max(xs[max(0,i-k+1):i+1]) for i in range(len(xs))]
\end{lstlisting}

A static analyzer run on each sample yields a set of inferred
post-conditions. Aggregating across all 20 samples, two post-conditions
survive in $\geq 90\%$ of samples: \texttt{len(out) <= len(xs)} and
\texttt{all(out[i] >= min(xs) for i)}. One post-condition, \texttt{len(out)
== len(xs) - k + 1}, survives in $14/20$ samples (the rest, like Sample
C, return a longer list). And one post-condition---\texttt{out[i] ==
max(xs[i:i+k])}---is fragile: it holds for Samples A and B but not C,
because C uses a left-aligned window.

The crucial observation: the \emph{disagreement} between samples is the
signal. The prompt under-specified what to do at the boundary. Samples
encode different boundary conventions; the analyzer makes those
differences legible as varying invariants. A single-sample pipeline sees
only one convention and treats it as ground truth. A distributional
pipeline sees the disagreement, can ask the user to disambiguate, can
use the dominant convention as a soft prior, and can flag the minority
samples as candidates for filtering.

%% --- 3. Approach --------------------------------------------------------
\section{Distributional Program Analysis}
\label{sec:approach}

\subsection{Setup}
Let $\phi$ be a specification (natural-language prompt and/or a partial
test suite). Let $M$ be an LLM. Sampling $N$ times from $M$ conditioned
on $\phi$ yields programs $p_1,\dots,p_N$. Let $A$ be an abstract
interpreter that, given a program $p$, returns a set $\mathcal{I}(p)$ of
\emph{candidate invariants}---boolean predicates over input/output pairs
or over internal state at named program points.

\subsection{Survival as posterior}
Let $\mathcal{I}^* = \bigcup_{i=1}^N \mathcal{I}(p_i)$ be the union of
candidate invariants observed across all samples. For each $\iota \in
\mathcal{I}^*$, define its \emph{survival rate} as
\[
s(\iota) = \frac{1}{N}\sum_{i=1}^N \mathbf{1}[\,A \text{ infers } \iota \text{ for } p_i\,].
\]
We use $s(\iota)$ as a posterior estimator: under the assumption that
the LLM samples concentrate around the intended program but with
independent decoding noise, an invariant of the intended program will
appear in a high fraction of samples; an invariant induced by noise will
appear in only a few. Formally, $s(\iota)$ is a Monte-Carlo estimate of
$\Pr_{p \samp M(\cdot|\phi)}[\,p \models \iota\,]$, which under the
above noise model approximates $\Pr[\,\iota \text{ holds of intended
program}\,]$.

\subsection{Decisions DPA enables}
The posterior $s(\cdot)$ is useful in three ways.
\begin{enumerate}
\item \emph{Sample filtering.} Reject samples that violate any invariant
with $s(\iota) \geq \tau$ for some threshold $\tau$ (we use
$\tau = 0.85$ in the pilot). Samples that disagree with the dominant
mode are likely off-distribution.
\item \emph{Specification refinement.} Surface the highest-survival
invariants to the developer as \emph{discovered} pre/post-conditions
that should be added to the specification, cleanly separating
``decoding noise'' from ``the model has internalized something the
prompt failed to say.''
\item \emph{Disagreement triage.} Invariants with mid-range survival
($0.4 < s < 0.7$) are typically under-specified design choices
(boundary conventions, sentinel values, error-handling). Surfacing them
turns a one-shot generation pipeline into a structured
clarification dialogue.
\end{enumerate}

\subsection{Pipeline}
Our prototype, on Python, performs five steps. (1) Sample $N=20$
programs from the LLM at temperature 0.8 with varied prompt
paraphrases. (2) Run an off-the-shelf interval and predicate-abstract
interpreter\footnote{We use a public Python static analysis stack; the
exact tool will be released with the supplementary appendix.} on each
sample, extracting candidate invariants at function entry, exit, and
each loop head. (3) Canonicalize invariants modulo $\alpha$-renaming
and constant folding. (4) Aggregate survival counts. (5) Apply the three
DPA decisions above on top of any existing downstream pipeline (test
selection, voting, repair).

\subsection{Cost model}
Naively, DPA is $N\times$ more expensive than analyzing one sample.
Crucially, however, DPA is intended to \emph{replace}, not supplement,
self-consistency-style aggregation~\cite{wang2023selfconsistency}, which
already requires $N$ samples. The marginal cost of DPA over
self-consistency is the abstract interpreter, which is small relative
to LLM inference (typically $<5\%$ in the pilot). The right
comparison is therefore not ``$N\times$ more compute'' but ``the same
compute with a sharper aggregator.''

%% --- 4. Pilot ------------------------------------------------------------
\section{Pilot Study}
\label{sec:pilot}

\paragraph{Setup.} We sampled $N=20$ programs per task from a frontier
LLM on the 164 tasks of HumanEval+~\cite{liu2023humanevalplus}, using
temperature 0.8 and three prompt paraphrases. For each task we ran the
analyzer, computed survival rates, and compared to a hand-extracted set
of reference invariants from the canonical solution.

\paragraph{Result 1: survival aligns with reference invariants.}
Pearson correlation between $s(\iota)$ and a binary indicator
``$\iota$ appears in the reference solution'' was $r = 0.78$ (p$<$0.001).
Among invariants with $s(\iota) \geq 0.85$, $82\%$ are present in the
reference solution. Among invariants with $s(\iota) \leq 0.30$, only
$11\%$ are present. The estimator is informative on both ends.

\paragraph{Result 2: filtering improves pass@1 at fixed compute.}
Table~\ref{tab:pilot} compares three aggregators applied to the same
20-sample budget: pick-first (random sample), self-consistency
(majority vote on test outputs), and DPA-filter (drop samples violating
any $s(\iota)\geq 0.85$ invariant, then majority vote on the rest).

\begin{table}[t]
\centering
\caption{HumanEval+ pilot. All conditions use the same 20 LLM samples
per task. Numbers are pass@1 rates over 164 tasks.}
\label{tab:pilot}
\begin{tabular}{@{}lc@{}}
\toprule
Aggregator & pass@1 \\
\midrule
Pick-first              & 64.6\% \\
Self-consistency~\cite{wang2023selfconsistency} & 71.3\% \\
DPA-filter (ours)       & \textbf{75.0\%} \\
\bottomrule
\end{tabular}
\end{table}

The 3.7-point absolute lift over self-consistency at \emph{matched
sample budget} is the central pilot finding. The lift is concentrated
on tasks with under-specified boundary behavior, exactly where
self-consistency's symptom-level voting fails: many wrong samples can
agree on a wrong output for a corner case, while the analyzer-derived
invariants disagree on the structure of the function and break the tie.

\paragraph{Qualitative observation.} On 18 tasks where DPA solved cases
self-consistency missed, manual inspection found that the surviving
invariants captured precisely the boundary convention the prompt left
ambiguous (off-by-one in indexing, empty-list return value,
exception-on-bad-input vs.\ silent default). The DPA-flagged disagreements
are, in nearly all cases, real specification gaps a developer would
want to know about.

%% --- 5. Evaluation plan & open questions --------------------------------
\section{Evaluation Plan and Open Questions}
\label{sec:plan}

\paragraph{Q1: Does the lift transfer beyond HumanEval+?}
We will replicate on MBPP+~\cite{liu2023humanevalplus},
LiveCodeBench~\cite{jain2024livecodebench} (which is harder to
contaminate), and a held-out internal corpus. Cross-benchmark
replication tests whether DPA exploits a general property of LLM code
distributions or a HumanEval-specific artifact.

\paragraph{Q2: Is DPA a substitute or a complement to capability?}
We will sweep three model classes (a frontier closed model, a strong
open-weight model, and a smaller open-weight model). We hypothesize
that DPA's lift is largest on the smaller model: weaker models produce
more decoding noise, leaving more variation for the analyzer to filter.

\paragraph{Q3: Does the framing extend beyond functional correctness?}
The pilot uses post-conditions over inputs/outputs. The same machinery
applies to security invariants (taint flow, no-deref-after-free) and
performance invariants (loop bounds, asymptotic complexity). We will
report a small case study on a taint-tracking analyzer.

\paragraph{Open question: contamination and prior collapse.}
If the LLM has memorized the canonical solution, $N$ samples may all be
near-copies of one program; survival rates would be uninformative
because the prior collapses. We will measure prior \emph{entropy}
(structural diversity across samples) per task and report DPA
performance as a function of entropy, to characterize the regime in
which the framing is meaningful versus degenerate.

\paragraph{Open question: false consensus.}
LLMs may share systematic biases inherited from pretraining
(e.g.\ Pythonic conventions that are unidiomatic in a specific domain).
A high-survival invariant could therefore reflect \emph{LLM consensus}
rather than \emph{intended-program truth}. We propose a control:
compare survival rates from a single model family to survival rates
across heterogeneous models; an invariant that survives \emph{across
families} is more credible.

%% --- 6. Threats ----------------------------------------------------------
\section{Threats to Validity}
\label{sec:threats}

\noindent\emph{Internal.} Survival rate is a Monte-Carlo estimate; for
small $N$, sampling variance is non-trivial. The pilot uses $N=20$,
which is on the small end; we report bootstrap confidence intervals in
the supplementary material. \emph{External.} The analyzer is one
abstract-interpretation stack; results may depend on which invariants
the underlying analyzer is capable of expressing. We mitigate by
ablating with a simpler interval analyzer in the planned full study.
\emph{Construct.} ``Reference invariants'' are extracted by the authors
and risk subjective bias; we will release the extraction protocol and
have a second annotator independently re-extract a subsample.
\emph{External.} Benchmark contamination is a known risk for
HumanEval-family benchmarks; the LiveCodeBench replication is the
primary mitigation.

%% --- 7. Related work ------------------------------------------------------
\section{Related Work}
\label{sec:related}

\paragraph{LLM aggregation.} Self-consistency~\cite{wang2023selfconsistency}
and majority-vote-style aggregators are the established way to exploit
multiple samples for a single task. They aggregate at the
\emph{output} level (which answer is most common). DPA aggregates at
the \emph{program-property} level, which is strictly richer:
self-consistency cannot distinguish two programs that produce the same
test output but differ in their loop bounds, error handling, or
aliasing behavior. AlphaCode-style filtering~\cite{li2022alphacode}
likewise filters by test agreement, not by inferred property.

\paragraph{Dynamic invariant inference.} Daikon~\cite{ernst2007daikon}
infers likely invariants by observing many \emph{executions} of one
program. DPA inverts this: we infer invariants by observing one
abstract execution of \emph{many} programs. The two are complementary;
a hybrid pipeline that runs Daikon on each LLM sample and then
aggregates is a natural extension.

\paragraph{Probabilistic programming.} The framing of the LLM as a prior
and the analyzer as a likelihood is closely analogous to probabilistic
programming~\cite{vandemeent2018ppintro}. We do not, in this paper,
take a Bayesian view formally; the survival rate is a frequentist
estimator. The connection is suggestive and motivates the formalism we
plan to develop in a follow-on full paper.

\paragraph{Static analysis on neural code.} Work on analyzing
LLM-generated code typically pairs an analyzer with a single sample
\cite{pearce2022copilotsec,sandoval2023llmsec}. We are unaware of prior
work that uses cross-sample variation as the central analytic signal.
SPoC~\cite{kulal2019spoc} and constrained-decoding work
\cite{poesia2022synchromesh} also use multiple samples but to filter
syntactically; DPA filters semantically through inferred properties.

%% --- 8. Conclusion --------------------------------------------------------
\section{Conclusion}
\label{sec:conclusion}

We have argued that an LLM, conditioned on a specification, is a
sampler over the under-determined program distribution that
specification permits. Existing pipelines treat that sampler as a
nuisance; we have proposed treating it as evidence. \textbf{Survival of
an inferred invariant across samples is a posterior estimator for the
properties of the latent intended program.} A small pilot supports
the claim. If the framing holds at scale, the implication for
LLM-driven software engineering is that the next gain will come not
from sampling more programs or running heavier analyses, but from
\emph{aligning} sampling-based generation with property-based analysis
through a Bayesian interface.

\balance

%% --- References ------------------------------------------------------------
\bibliographystyle{ACM-Reference-Format}
\bibliography{ref}

\end{document}
